\section{Refinement and consistency}\label{app:execution}

\subsection{The refinement cost guarantee}\label{app:refinement-guarantees}\label{sec:refinement-guarantee}

We prove Lemma~\ref{lem:refinement-cost}. The successful central choice
was analysed in \eqref{eq:certified-anchor}; the remaining case chooses
an anchor in the middle half of the coordinate's interval.

\begin{proof}[Proof of Lemma~\ref{lem:refinement-cost}]
The in-node DAG bounds ensure that the number of other coordinates
comparable with $x$ is at most $B$ from \eqref{eq:comparable-bound}.
Suppose the routine uses the middle-half interval. Either comparison
certification failed, or a structural refinement's call to
\texttt{CentralInterval} failed. In both cases,
Lemma~\ref{lem:certificate-success} gives $W_x<w_*=2^{-2kB}$.
Write $L=\log(1/W_x)>2kB$. The coordinate-density bound
\eqref{eq:fiber-lower} and the refinement identity give
\[
g=\log f_x(c)\ge L-2B.
\]
Since the rounded interval has length $W_x/2$,
Lemma~\ref{lem:rounding} gives $m<L+1$. Consequently
\[
g\ge L-2B>(1-1/k)L>(1-1/k)m-4.
\]
The central choice satisfies the same final bound, by
\eqref{eq:certified-anchor}. Thus in every case
\begin{equation}\label{eq:tariff}
g\ge(1-1/k)m-4.
\end{equation}
In particular $g+4\ge0$.
If the refinement follows failed comparison certification, it uses the
middle half directly and has
\begin{equation}\label{eq:failure-gain}
g\ge L-2B>2B(k-1)>k.
\end{equation}
Both choices return an interior dyadic anchor and terminate.

Combining \eqref{eq:tariff} with the child cost
\eqref{eq:anchor-cost} gives
\[
\ell(c)\le\frac{\gamma}{1-1/k}(g+4)+\log k+1.
\]
For $k\ge64$, the estimate
$-\log(1-1/k)\le1/((k-1)\ln2)\le3/(2k)$ gives
$\gamma/(1-1/k)\le1+3/k$. Set
\[
a=1+\frac3k,\qquad b=4a+\log k+1.
\]
Because $g+4\ge0$, increasing its coefficient to $a$ preserves the
inequality. Thus $\ell(c)\le ag+b$, including when $g$ is negative.
The coefficients have $a=1+O(1/k)$ and $b=O(\log k)$, with
$a\ge1$ and $1\le b\le\log k+6<k+1$. For $k\ge64$, they also
satisfy $a\le3/2$ and $b\le k/4$, since $\log k+6\le k/4$.
\end{proof}

\subsection{Consistency and termination}

We first justify admission to an existing child, then verify consistency
through an entire comparison call. Assignments with equal free
coordinates, or with a free coordinate equal to an anchor, form a finite
union of sets of volume zero. They are excluded when reading a total
order from the state.

\begin{proof}[Proof of Lemma~\ref{lem:admission}]
Consider first an existing child at $c$, and denote its label set by $S$.
Choose a feasible assignment in the parent section $X_x=c$, with the other
free coordinates pairwise distinct and different from every child anchor.
Such an assignment exists by Lemma~\ref{lem:interior}. Also choose feasible
assignments in the other nodes.

There is a positive distance from $c$ to every other child anchor and
every free parent coordinate in this assignment. All incident inequalities at $x$ are strict. Thus, with
everything else fixed, $x$ can be moved slightly to either side of $c$
while keeping the old state feasible. The resulting two represented orders
place $x$ immediately before or immediately after the entire block $S$.
Both orders respect the transcript. For each $y\in S$, one has $x<y$ and
the other has $y<x$; neither relation can therefore be implied by the
transcript.

Now take any order represented after the refinement. Deleting $x$ from the
child leaves an assignment of its old coordinates and hence an old child
order. Apply the preceding perturbation using this assignment's parent
coordinates and old child coordinates. It gives two feasible old orders
that agree with the new one on every label except $x$. The new order
inserts $x$ somewhere into their child block. No
answered comparison joins $x$ to a label in $S$, so this insertion violates
none of them. Relative to every label outside $S\cup\{x\}$, the whole
block occupies the old anchor position. The item $x$ is consequently on
the same side of each outside label as in either of the two feasible old
orders above. Comparisons between other labels are unchanged. The new
order therefore respects the entire transcript.

If there was no child at $c$, the same argument uses an empty block: the
new singleton child places $x$ where its fixed parent coordinate stood.
Finally, the remaining parent region is a positive-volume section, and
the child receives an independent coordinate in $(0,1)$. All node regions
remain positive.
\end{proof}


\begin{proof}[Proof of Proposition~\ref{prop:legality}]
Initially every node has positive volume, every order is represented,
and the strict in-node DAG bounds hold. We induct over completed calls
and, within each call, over its changes to the state.

Before answering, a refinement deletes one vertex and adds an isolated
vertex to a child, so the strict DAG bounds continue to hold. The anchor
routine returns an interior position, as proved in
Appendix~\ref{app:refinement-guarantees}. Lemma~\ref{lem:admission}
preserves positive node volumes and consistency with the transcript.
Normalization preserves the represented region by
Lemma~\ref{lem:interior}.

If the queried labels lie in different child subtrees, their anchors
determine the answer in every represented order. Otherwise, the
acceptance test in Algorithm~\ref{alg:compare} and
Lemma~\ref{lem:estimation} ensure that the chosen answer retains at least
$1/2-\eta>0$ of its node's volume. Recording its arc or interval bound
therefore preserves positive volume and makes the answer true in every
represented order. A new arc cannot create a cycle, since that would
make the region empty. Normalization preserves this updated region.
The transcript is then extended by the chosen answer.

During the subsequent repair, the temporary DAG bounds follow from
Lemma~\ref{lem:restoration}. Each move again preserves positive volume
and the transcript by Lemma~\ref{lem:admission}. This applies to every
vertex of a selected path, including those whose incident arcs were
removed earlier in the batch. The full batch is completed, and
Lemma~\ref{lem:restoration} restores the strict bounds for the next call.

The finite certificate calculations and the explicit rounding rule
terminate. As shown in Section~\ref{sec:repairs}, before answering each
queried label can be refined only once: a second move would have gain
zero, contradicting the failed-certification gain guarantee. Thus there
are at most two such moves. A repair makes at most $n$ further moves,
so every call terminates with at most $n+2$ refinements.
All analytic guarantees used here apply to arbitrary positive-volume
regions with the stated DAG bounds, independently of the interaction
history.
\end{proof}
